\documentclass[border = 1mm]{standalone}

% --
\usepackage[utf8]{inputenc}
\usepackage{polyglossia}
\setmainlanguage{spanish}

\usepackage{amsmath}

% --
\usepackage{tikz}
\usetikzlibrary{
	math,
	arrows.meta,
}

\begin{document}
\begin{tikzpicture}[
	line cap = round,
	line join = round,
]
	
	%. tamaño total del dibujo
	\tikzmath{\L = 6.5;}
	
	%. márgenes
	\draw[gray!50] (0,0) rectangle (\L,\L);
	%. máscara de recorte
	\clip (0,0) rectangle (\L,\L);

	% -- dibujo principal:
	\begin{scope}[	
		%. definir el oirigen en el centro del dibujo
		shift = {(0.5*\L,0.5*\L)},
	]
	
	% A. ecuaciones
	\node at (0,2.35) {
	$\displaystyle
		\mathcal{F}\big[ \operatorname{rect}(x/L) \big] =
		L \operatorname{sinc}(L u)
	$
	};
	
	% B. funciones
	
	% rect(x)
	\begin{scope}[
		shift = {(-1.8,0.15)},
		scale = 0.85,
	]
		
		% ejes cartesianos
		\draw[gray!50]
			(0,-0.2) to (0,1.2) node[above, black] {\small $f(x)$}
			(-1.2,0) to (1.2,0) node[right, black] {\small $x$};
		
		% función
		\tikzmath{ \L = 0.75; \H = 0.75; }
		
		\draw[line width = 0.5pt]
			(-1.2,0) to (-0.5*\L,0)
			(-0.5*\L,\H) to (0.5*\L,\H)
			(0.5*\L,0) to (1.2,0);
		\draw[dashed]
			(-0.5*\L,0) to ++(0,\H)
			(0.5*\L,0) to ++(0,\H);
		
		\foreach \p in {(-0.5*\L,0), (-0.5*\L,\H), (0.5*\L,0), (0.5*\L,\H)}
			\draw[line width = 0.5pt, fill = white] \p circle (0.03);
			
		\foreach \p in {(-0.5*\L,0.5*\H), (0.5*\L,0.5*\H)}
			\draw[line width = 0.5pt, fill] \p circle (0.025);
	
		% --
		\draw[arrows = {Bar-Bar}]
			(-0.5*\L,-0.4) to
				node[below] {\footnotesize $L$}
			(0.5*\L,-0.4);
		
	\end{scope}
	
	%. sinc(x)
	\begin{scope}[
		shift = {(+1.3,0.15)},
		scale = 0.85,
	]
		
		% ejes cartesianos
		\draw[gray!50]
			(0,-0.2) to (0,1.2) node[above, black] {\small $F(u)$}
			(-1.5,0) to (1.5,0) node[right, black] {\small $u$};
		
		% ticks
		\foreach \i in {-5,-4,...,5}
			\draw[gray!50] (\i/3.5,0.08) to ++(0,-0.16);
		
		% función
		\draw[line width = 0.6pt, smooth]
			plot[domain = -1.5:-0.005, samples = 80]
			(\x, {sin(3.5*pi*\x r)/(3.5*pi*\x)})
			--
			plot[domain = 0.005:1.5, samples = 80]
			(\x, {sin(3.5*pi*\x r)/(3.5*pi*\x)});
			
		\draw[line width = 0.5pt, fill] (0,1) circle (0.025);
		
		\node[below = 1.5mm, xshift = -1.5mm] at (-1/4,0)
			{\footnotesize $\displaystyle -\frac{\pi}{L}$};
		\node[below = 1.5mm] at (1/4,0)
			{\footnotesize $\displaystyle \frac{\pi}{L}$};
		
	\end{scope}
	
	% C. Flecha
	\draw[
		semithick, line cap = butt,
		gray!75,
		double, double distance = 0.6mm,
		arrows = {Implies-Implies},
	]
		(-0.7,0.8) to
			node[above = 0.5mm, black] {\small $\mathcal{F}$}
		++(0.85,0);
	
	% D. definiciones
	\node at (-1.75,-1.65) {
	\footnotesize
	$\displaystyle
	\begin{cases}
		0, & |x|>\frac{L}{2} \\[1.5mm]
		\frac{1}{2}, & |x| = \frac{L}{2} \\[1.5mm]
		1, & |x|< \frac{L}{2}.
	\end{cases}
	$
	};

	\node at (+1.4,-1.65) {
	\footnotesize
	$\displaystyle
	\begin{cases}
		\dfrac{\sin(L u)}{L u}, & u\neq 0 \\[3mm]
		1, & u = 0.
	\end{cases}
	$
	};

	\end{scope}
	% --

\end{tikzpicture}
\end{document}